\documentclass[11pt,a4paper]{article}
\usepackage[T2A]{fontenc}
\usepackage[utf8]{inputenc}
\usepackage[russian]{babel}
\usepackage{booktabs}
\usepackage{geometry}
\usepackage{amsmath}
\geometry{margin=2.2cm}

\title{Изменение формы REM-профилей вокруг сейсмособытий\\
после исключения явных артефактов}
\author{}
\date{13 июля 2026}

\begin{document}
\maketitle

\section{Цель}
Проверить, меняется ли форма суточного REM-профиля при приближении к
сейсмособытию. Основные признаки --- расстояния L1 и L2 каждого дня до
поточечного среднего профиля первых трёх дней окна. Дополнительно оценивались
среднее и размах профиля.

\section{Данные и предобработка}
Использованы 8-дневные профили с 24 точками в сутки, \texttt{overlap=0.50},
\texttt{rem\_stage=2}. Каждый полный 8-дневный вектор нормировался в
$[-1,1]$ своей min--max нормировкой. Следовательно, результаты описывают
изменение \emph{формы} профиля, но не абсолютного уровня REM.

Из анализа исключены явные артефакты:
\begin{itemize}
  \item строка 0, день 1;
  \item строки 4 и 5, день 3;
  \item строка 22, день 4.
\end{itemize}
Исключённый день не участвует ни в расчёте своей дневной метрики, ни в
референсном профиле. В частности, если артефакт попадает в первые три дня,
референс усредняется только по оставшимся валидным дням.

Для каждого окна референсный профиль равен
\[
  r_i = \frac{1}{|B|}\sum_{d\in B} y_{d,i}, \qquad
  B=\{1,2,3\}\setminus\{\text{артефактные дни}\}.
\]
Для каждого дня $d$ рассчитывались
\[
  L1_d=\sum_i |y_{d,i}-r_i|,
  \qquad
  L2_d=\sqrt{\sum_i(y_{d,i}-r_i)^2}.
\]
Сравнивались средние L1/L2 на позднем участке (дни 6--8) и раннем участке
(дни 1--3; только валидные дни).

\section{Статистический протокол}
Крысы, зарегистрированные на одну дату события, не рассматриваются как
независимые наблюдения. Поэтому сначала признаки усреднялись внутри даты
события.

Изначальное сравнение ``дни 6--8 минус дни 1--3'' для L1/L2 нельзя
использовать как подтверждающий тест: дни 1--3 участвуют в построении
референсного профиля и поэтому по конструкции имеют меньшую дистанцию до
него. Это нарушает обменность ранних и поздних наблюдений, необходимую для
парного sign-flip test. Такие значения $p$ могут использоваться только как
описательные.

В качестве менее смещённой проверки выполнено сравнение дней 6--8 с днями
4--5. Обе группы не входят в построение референса. Для разности
``дни 6--8 минус дни 4--5'' приведён точный sign-flip test на уникальных
датах. Он всё ещё предполагает независимость дат и не устраняет возможную
временную автокорреляцию; поэтому является исследовательской, а не
подтверждающей проверкой.

\section{Основной результат}
В окнах \texttt{before} было 18 профилей на 11 уникальных дат, в
\texttt{after\_reversed} --- 15 профилей на 8 дат. В \texttt{after\_reversed}
индекс дня направлен от более далёкого дня после события к более близкому:
день 8 соответствует ближайшему дню после события.

\begin{table}[ht]
\centering
\caption{Описательные расстояния до среднего профиля первых трёх дней после
очистки от артефактов. Эти значения не являются основой для $p$-теста,
поскольку ранние дни входят в референс.}
\begin{tabular}{llrrr}
\toprule
Окно & Метрика & Дни 1--3 & Дни 6--8 & Разность \\
\midrule
before & L1 & 4.496 & 7.859 & +3.363 \\
before & L2 & 1.120 & 1.939 & +0.820 \\
after\_reversed & L1 & 5.243 & 7.402 & +2.159 \\
after\_reversed & L2 & 1.267 & 1.825 & +0.558 \\
\bottomrule
\end{tabular}
\end{table}

\begin{table}[ht]
\centering
\caption{Нециклическая исследовательская проверка: дни 6--8 против дней
4--5 для расстояний до референса первых трёх дней.}
\begin{tabular}{llrrrrr}
\toprule
Окно & Часть профиля & Метрика & Разность & Знак & $p_{\rm one}$ & $p_{\rm two}$ \\
\midrule
before & все точки & L1 & +1.275 & 5/11 & 0.1133 & 0.2266 \\
before & все точки & L2 & +0.272 & 5/11 & 0.1489 & 0.2979 \\
before & первая половина & L1 & +0.208 & 5/11 & 0.3457 & 0.6914 \\
before & первая половина & L2 & +0.022 & 4/11 & 0.4443 & 0.8887 \\
before & вторая половина & L1 & +1.068 & 8/11 & 0.0449 & 0.0898 \\
before & вторая половина & L2 & +0.306 & 8/11 & 0.0684 & 0.1367 \\
after\_reversed & все точки & L1 & -0.254 & 3/8 & 0.5898 & 0.8281 \\
after\_reversed & все точки & L2 & +0.021 & 3/8 & 0.4727 & 0.9453 \\
\bottomrule
\end{tabular}
\end{table}

При раздельном описательном расчёте по половинам профиля рост L1/L2 от
дней 1--3 к дням 6--8 виден в обеих половинах окна \texttt{before}.
Однако это наблюдение сохраняет проблему референса и не должно
интерпретироваться по соответствующим $p$-значениям.

\section{Сравнение внутриблочных расстояний без референса}
Чтобы полностью исключить зависимость от эталонного профиля, для каждого
окна рассчитана средняя попарная дистанция внутри раннего и позднего блоков:
\[
  D_q(B)=\frac{1}{\binom{|B|}{2}}
  \sum_{i<j,\;i,j\in B}\lVert y_i-y_j\rVert_q,
  \qquad B\in\{\{1,2,3\},\{6,7,8\}\}.
\]
Для L1 используется $q=1$, для евклидовой метрики --- $q=2$. При наличии
артефакта он исключается из соответствующего блока; например, для строки 0
ранний блок содержит только дни 2--3. Эта статистика отвечает на другой
вопрос: насколько дни \emph{внутри} позднего блока похожи друг на друга
относительно раннего блока.

\begin{table}[ht]
\centering
\caption{Средние попарные расстояния внутри блоков; тестируется разность
$D(6\text{--}8)-D(1\text{--}3)$ после усреднения по крысам одной даты.}
\begin{tabular}{llrrrrr}
\toprule
Окно & Часть & Метрика & $D(1\text{--}3)$ & $D(6\text{--}8)$ &
$p_{\rm one}$ & $p_{\rm two}$ \\
\midrule
before & все точки & L1 & 7.990 & 8.712 & 0.1875 & 0.3750 \\
before & все точки & L2 & 1.984 & 2.196 & 0.1323 & 0.2646 \\
before & первая половина & L1 & 3.891 & 3.945 & 0.4722 & 0.9443 \\
before & первая половина & L2 & 1.342 & 1.322 & 0.5430 & 0.9150 \\
before & вторая половина & L1 & 4.099 & 4.767 & 0.1265 & 0.2529 \\
before & вторая половина & L2 & 1.365 & 1.604 & 0.0942 & 0.1885 \\
after\_reversed & все точки & L1 & 9.015 & 7.155 & 0.9766 & 0.0547 \\
after\_reversed & все точки & L2 & 2.192 & 1.790 & 0.9883 & 0.0312 \\
\bottomrule
\end{tabular}
\end{table}

Для \texttt{before} надёжного отличия внутриблочных расстояний не найдено.
Для \texttt{after\_reversed} L2 внутри дней 6--8 меньше, чем внутри дней
1--3 ($p_{\rm two}=0.0312$); это противоположно заранее направленной
гипотезе о росте вариабельности и требует осторожной репликации.

\section{Проверка устойчивости к параметрам профиля}
Доступный кэш позволяет проверить 9 сочетаний числа точек в сутки
$N\in\{12,24,48\}$ и overlap $\in\{0,0.25,0.50\}$; во всех случаях
\texttt{rem\_stage=2}. В этих кэшах имеется 17 профилей \texttt{before}
на 11 дат и 14 профилей \texttt{after\_reversed} на 8 дат. Те же исключения
артефактов применялись по идентификатору крысы, дате и направлению окна.

Во всех девяти конфигурациях описательно наблюдался рост L1/L2 от дней 1--3
к дням 6--8. Это подтверждает устойчивость \emph{описательного} паттерна к
разрешению и overlap, но не делает исходные $p$-значения валидными: во всех
конфигурациях сохраняется та же зависимость референса от первых трёх дней.
Абсолютные L1/L2 закономерно меняются с числом точек, поэтому сопоставимы
только направление и стандартизированный размер эффекта.

\section{Выводы и ограничения}
\begin{enumerate}
  \item L1/L2 на поздних днях описательно выше, чем на первых трёх днях,
  во всех проверенных конфигурациях. Этот результат не является
  подтверждающим из-за включения ранних дней в референс.
  \item В нециклическом сравнении дней 6--8 с днями 4--5 убедительный
  двусторонний эффект отсутствует. Есть только исследовательский сигнал L1
  во второй половине профиля окна \texttt{before}
  ($p_{\rm one}=0.0449$, $p_{\rm two}=0.0898$).
  \item В сравнении попарных расстояний внутри блоков без референса эффект
  для \texttt{before} не подтверждается. В \texttt{after\_reversed} L2
  показывает меньшую внутриблочную вариабельность на днях 6--8
  ($p_{\rm two}=0.0312$), что требует отдельной проверки.
  \item До биологической интерпретации ``день/ночь'' необходимо проверить,
  что первая и вторая половины точек действительно соответствуют нужным
  фазам светового цикла, а не только половинам массива.
  \item Для подтверждающего анализа нужны заранее заданные артефактные
  критерии, независимый контрольный период или псевдодаты и перестановочный
  тест, в котором при каждой перестановке референс пересчитывается заново.
\end{enumerate}

\end{document}
